% RESULTADO_RECUPERADO desde II REV09. Véase MAPA_PROCEDENCIA.json.
% Selección de líneas y cambios mecánicos explícitos; ninguna prueba nueva.
\subsection{Normalización del estado y significado del Hamiltoniano modular}
El segundo observable se explicita mediante el estado que lo define.
Se mantiene la descomposición tensorial anterior: dieciocho qutrits
por sector, cada uno con operador número $\operatorname{diag}(0,1,2)$;
$N_m$ es la suma de esos operadores y los sectores actúan en factores
distintos. Sea
\[
 z_m=1+e^{-m\Delta_{\rm mod}}+e^{-2m\Delta_{\rm mod}},
 \qquad m\in\{90,120\}.
\]
La familia de estados de la fuente tiene la realización producto
\[
 \rho_A=Z^{-1}e^{-\Delta_{\rm mod}(90N_{90}+120N_{120})},
 \qquad Z=z_{90}^{18}z_{120}^{18}.
\]
La conmutatividad de los operadores número y la factorización tensorial
prueban la expresión de $Z$. El estado es de rango completo para
$\Delta_{\rm mod}$ real y finito, y su logaritmo espectral da
\[
 K_A=-\log\rho_A=(\log Z)I+\Delta_{\rm mod}D.
\]
Para $\Delta_{\rm mod}\ne0$, $K_A$ pesa de forma distinta las excitaciones
de los dos sectores. La lectura conjunta con $N$ recupera las dos ocupaciones
por la ecuación~\eqref{v:ant:eq:recover}.

La necesidad de la lectura conjunta se ve en dos ejemplos complementarios.
Los estados de ocupación $(1,0)$ y $(0,1)$ tienen igual total $N=1$
y diferente peso modular. A la inversa, $(4,0)$ y $(0,3)$ tienen
el mismo $D=360$ y diferentes totales. La pareja $(N,D)$ separa ambos
casos. Esta recuperación se refiere a ocupaciones sectoriales, no a
cada vector microscópico dentro de una degeneración de esas ocupaciones.

Cuando $\Delta_{\rm mod}=0$, el estado es proporcional a la identidad
y $K_A$ es escalar: esa lectura deja de distinguir los sectores.
Para $\Delta_{\rm mod}\ne0$, la normalización conocida permite recuperar
$D$ y componerlo con el área. Así queda definido el papel del Hamiltoniano
modular: caracteriza el estado reducido y su distribución de información.
Su inclusión en el artículo está motivada por esa función matemática,
mientras el Hamiltoniano de evolución pertenece al desarrollo dinámico
de la acción.
